\chapter{Réalisation et fonctionnalités}

\minitoc

% ------------------------------------------------------------------------------------
% UTILISATION DE CE CHAPITRE
% Chaque module ou fonctionnalité principale du projet dispose d'une \section selon
% le même modèle en deux parties : « Fonctionnalités réalisées » (réalisations, choix
% de conception et contraintes), puis « Captures d'écran » (interface ou résultats).
% Copier le bloc \section ci-dessous autant de fois que nécessaire, puis supprimer
% ce commentaire et les deux sections d'exemple.
% ------------------------------------------------------------------------------------

\section{[Nom de la fonctionnalité / du module]}

\subsection{Fonctionnalités réalisées}
% TODO: présenter les réalisations, les choix de conception, les cas particuliers et les contraintes.
[Décrivez les réalisations associées à cette fonctionnalité.]

\subsection{Captures d'écran}
% TODO: ajouter une figure par capture d'écran, par exemple :
% \begin{figure}[H]
%     \centering
%     \includegraphics[width=0.8\textwidth]{Chapter4/images/your-screenshot.png}
%     \caption{[Légende décrivant l'écran]}
% \end{figure}
[Insérez les captures d'écran ici.]

\section{[Nom de la fonctionnalité / du module]}

\subsection{Fonctionnalités réalisées}
[Décrivez les réalisations associées à cette fonctionnalité.]

\subsection{Captures d'écran}
[Insérez les captures d'écran ici.]

\section{Conclusion}
% TODO: conclure le chapitre en 2 à 3 phrases.
[Rédigez une brève conclusion du chapitre.]
